\chapter{Architektur}\label{kap:architektur}

\section{Komponenten und Datenfluss}

Das System besteht aus einem einzigen Docker-Container, in dem zwei Prozesse
laufen: der R-Scheduler und der Rust-Simulationsserver. Dazu kommen ein
Docker-Volume und ein Webserver.

\begin{figure}[ht]
\centering
\resizebox{\textwidth}{!}{\begin{tikzpicture}[
  font=\sffamily\footnotesize,
  box/.style={draw=ink3,fill=white,rounded corners=1.5pt,align=center,inner sep=5pt,minimum height=1.05cm},
  ext/.style={box,fill=surface2},
  arr/.style={-{Stealth[length=5pt]},draw=ink2,thick},
  lab/.style={font=\sffamily\scriptsize,color=ink3,align=center},
  grp/.style={draw=rule,dashed,rounded corners=3pt,inner sep=8pt},
]
\node[ext,text width=2.1cm] (api) {API-Football\\[1pt]\scriptsize(RapidAPI)};
\node[box,right=1.3cm of api,text width=2.6cm] (sch) {\textbf{R-Scheduler}\\[1pt]\scriptsize 11--23 Uhr, Takt 2 min};
\node[box,below=1.15cm of sch,text width=2.6cm] (rust) {\textbf{Rust-Server}\\[1pt]\scriptsize axum + rayon, :8080};
\node[box,right=1.6cm of sch,text width=2.4cm] (gen) {\textbf{Seitengenerator}\\[1pt]\scriptsize R + htmltools};
\node[ext,right=1.5cm of gen,text width=2.1cm,shape=cylinder,shape border rotate=90,aspect=0.18,minimum height=1.25cm] (vol) {Volume\\[1pt]\scriptsize\mbox{\texttt{fussball-site}}};
\node[box,below=2.8cm of vol,text width=2.1cm] (caddy) {\textbf{Caddy}\\[1pt]\scriptsize HTTPS, Cache};
\node[ext,below=2.8cm of gen,text width=2.4cm] (browser) {Browser\\[1pt]\scriptsize\mbox{\texttt{verlauf.js}}, Veraltet-Hinweis};
\node[ext,below=1.15cm of api,text width=2.1cm] (team) {TeamList-CSV\\[1pt]\scriptsize Start-ELO};

\begin{scope}[on background layer]
  \node[grp,fit=(sch)(rust)(gen),label={[lab,anchor=south west]north west:Container \texttt{fussball-scheduler}}] {};
\end{scope}

\draw[arr] (api) -- node[lab,above]{Spiele} (sch);
\draw[arr] ([xshift=-6pt]sch.south) -- node[lab,left]{JSON} ([xshift=-6pt]rust.north);
\draw[arr] ([xshift=6pt]rust.north) -- node[lab,right]{Matrix,\\Details} ([xshift=6pt]sch.south);
\draw[arr] (sch) -- node[lab,above]{Ergebnisse} (gen);
\draw[arr] (gen) -- node[lab,above]{22 Seiten} (vol);
\draw[arr] (vol) -- node[lab,right]{nur lesend} (caddy);
\draw[arr] (caddy) -- node[lab,above]{HTTPS} (browser);
\draw[arr,dashed] (team) -- (sch);
\end{tikzpicture}}
\caption{Komponenten zur Laufzeit. Der Container ist gestrichelt umrandet.}
\label{fig:architektur}
\end{figure}

\begin{enumerate}
  \item \textbf{API-Football} über RapidAPI liefert Spielpläne, Ergebnisse
        und Live-Status.
  \item Der \textbf{R-Scheduler} fragt die API im Zwei-Minuten-Takt ab,
        bereitet die Daten auf und ruft den Rust-Server.
  \item Der \textbf{Rust-Server} rechnet die Monte-Carlo-Simulation und die
        deterministischen Spieldetails. Er ist nur innerhalb des Containers
        erreichbar.
  \item Der \textbf{Seitengenerator} (R, \code{htmltools}) schreibt die 22
        Seiten in das Ausgabeverzeichnis. Die Ergebnisse übergibt der
        Scheduler im Speicher. Im Produktionspfad wird keine Ergebnisdatei
        geschrieben.
  \item Das \textbf{Docker-Volume} \code{fussball-site} teilen Container und
        Webserver. Der Container schreibt, Caddy liest.
  \item \textbf{Caddy} liefert die Seiten mit HTTPS, Kompression und
        passenden Cache-Regeln aus.
\end{enumerate}

\section{Leitentscheidungen}

Die Architektur folgt einer Reihe festgehaltener Entscheidungen
(\datei{docs/adr/}):

\begin{center}
\small
\begin{tabularx}{\textwidth}{@{}lX@{}}
\toprule
ADR & Entscheidung \\
\midrule
0001 & Statische Seiten statt einer gehosteten Shiny-App. \\
0002 & Spieldetails werden bei jedem Rendern deterministisch aus Startwerten und Ergebnissen berechnet und nicht gespeichert. Die Modellkonstanten leben nur im Rust-Server. \\
0003 & ELO-Eichung für Ligen ohne Historie. \\
0004 & Ein Tormodell je Wechselgemeinschaft. \\
0005 & Ligakonfiguration als Daten in einer Registry. \\
0006 & Abstiegskopplung der Regionalligen an die 3.~Liga, exakt ausgezählt. \\
0007 & Die TeamList ist ein gepflegtes Stammdatenblatt, der Saisonwechsel läuft in zwei Phasen. \\
0008 & Zweitvertretungs-Malus nach §\,55b Nr.~3.1 DFB-Spielordnung. \\
\bottomrule
\end{tabularx}
\end{center}

Eine frühere Aufteilung in Microservices mit Kubernetes wurde im Mai 2026
zugunsten des einen Containers zurückgenommen (Git-Tag
\code{pre-deployment-cleanup-2026-05-02}). Diagramme dieses Entwurfs in
\datei{docs/architecture/} sind als historisch markiert.
